% Generated draft: requires mathematical and source review.
\documentclass[11pt]{article}
\usepackage[margin=25mm]{geometry}
\usepackage{fontspec}
\setmainfont{Latin Modern Roman}
\usepackage{amsmath,amssymb,mathtools,unicode-math}
\setmathfont{Latin Modern Math}
\usepackage{enumitem,needspace,xurl,hyperref}
\setlength{\parindent}{0pt}
\setlength{\parskip}{5pt}
\setlength{\emergencystretch}{3em}
\newcommand{\N}{\mathbb{N}}
\newcommand{\Z}{\mathbb{Z}}
\newcommand{\Q}{\mathbb{Q}}
\newcommand{\R}{\mathbb{R}}
\newcommand{\C}{\mathbb{C}}
\newcommand{\F}{\mathbb{F}}
\newcommand{\A}{\mathbb{A}}
\newcommand{\PP}{\mathbb{P}}
\newcommand{\E}{\mathbb{E}}
\newcommand{\eps}{\varepsilon}
\newcommand{\dd}{\,\mathrm{d}}
\begin{document}
\section*{UnsolvedMath ID: 30004478}
\textbf{Source problem code:} catalogue entry
\subsection{Definitions and mathematical statement}
Let \(F\) be a number field, that is, a finite extension of \(\Q\), and let \(C/F\) be a smooth, geometrically irreducible projective curve of genus \(g(C)\geq 2\). Geometric irreducibility means that \(C\times_F\overline F\) is irreducible. Write \(C(F)\) for the set of \(F\)-rational points of \(C\). For a projective presentation \(C\subseteq\PP_F^n\), these are the projective coordinate tuples in \(\PP^n(F)\) satisfying the defining equations, with tuples identified under multiplication by a common nonzero element of \(F\).

\textbf{Algorithmic question.} Does there exist a single algorithm which, for every number field \(F\) and every such curve \(C/F\), takes exact equations defining \(C\), together with an effective finite presentation of \(F\) and its coefficients, halts, and outputs a finite list of distinct points \(P_1,\ldots,P_m\) such that
\[
C(F)=\{P_1,\ldots,P_m\}?
\]
Here \(m\geq 0\), and the output must be the empty list when \(C(F)=\varnothing\). An effective presentation of the number field and exact coordinates is the computational interpretation of the supplied phrase ``given equations''; the record does not prescribe a particular encoding. The input is promised to satisfy the stated geometric hypotheses.

\textbf{Additional question in the original statement.} The supplied original formulation also asks whether \(\#C(F)\) can be bounded above using suitable known invariants of \(C\). It does not specify which invariants are allowed, whether they must be computable from the input, or what effectiveness and uniformity are required of the bound. This additional target is retained here as an unresolved formulation qualification; it is absent from the supplied clean statement. The original wording also does not explicitly include geometric irreducibility, whereas the clean statement does.

The supplied literature assessment reports uniform cardinality bounds involving genus, field degree, and Jacobian rank, while distinguishing these from an algorithm producing the complete list. This is supplied dated curation, not an independently verified theorem statement or literature assessment.
\subsection{Short English statement}
The main question asks for one terminating algorithm that lists every rational point on any given smooth geometrically irreducible projective curve of genus at least \(2\) over a number field, including certifying an empty list when appropriate. The original record additionally asks for cardinality bounds in unspecified known invariants; that part needs further formulation. The supplied open label and literature assessment are unverified curation.
\subsection{Sources}
\begin{itemize}
\item Local UnsolvedMath registry, permanent ID \(30004478\), title \emph{Effective Enumeration of Rational Points on Curves}, source problem code \texttt{catalogue entry}, in the Oberwolfach Reports Open Problems collection. The supplied record gives proposed year \(2020\) and no named proposer.
\item Supplied source citation: \emph{Arithmetic Geometry (2021)}, Oberwolfach Reports, pp.\ 1023--1082, DOI \texttt{10.4171/owr/2020/20}. Original OWR report: \url{https://doi.org/10.4171/owr/2020/20}. The supplied background describes these reports as records of mathematical workshops and their open-problem discussions. The citation is reproduced as supplied; its year and exact problem location have not been independently checked.
\item Supplied literature source attributed to Dimitrov--Gao--Habegger: \url{https://annals.math.princeton.edu/2021/194-1/p04}. The record attributes strong uniform cardinality bounds depending on genus, field degree, and Jacobian rank to this work, and says these do not provide the requested complete-list algorithm. No precise theorem or hypotheses are supplied here.
\item Literature assessment dated exactly \texttt{2026-08-22}: the supplied status is \texttt{open}, and the assessment says that an effective complete-list algorithm for an arbitrary curve in the stated setting remains unknown. This is supplied dated curation derived from the OWR record, not independently verified. The record itself requests independent verification of time-sensitive status and cited claims.
\item The supplied statement status is \texttt{corrected\_verified}, with a registry assertion that the clean statement was checked against the official report and preserves its hypotheses, notation, and open target. That assertion is preserved as source metadata; this draft does not independently verify it.
\end{itemize}
\end{document}
